% Generated from paper002_supplement_model_order_v1.1.md by scripts/build_paper002_submission_tex.py.
% The Markdown and BibTeX files remain the synchronized content sources.
\documentclass[10pt]{article}
\usepackage[a4paper,margin=25mm]{geometry}
\usepackage[T1]{fontenc}
\usepackage[utf8]{inputenc}
\usepackage{amsmath,amssymb}
\usepackage{graphicx}
\usepackage{booktabs,tabularx,array}
\usepackage{enumitem}
\usepackage{fancyvrb}
\usepackage{caption}
\usepackage[numbers,sort&compress]{natbib}
\usepackage{xurl}
\usepackage[colorlinks=true,allcolors=blue]{hyperref}
\setlength{\parindent}{0pt}
\setlength{\parskip}{0.55em}
\renewcommand{\arraystretch}{1.12}
\graphicspath{{figures/}}
\renewcommand{\thesection}{S\arabic{section}}
\title{Supplement: Failure-Conditioned Model-Order Expansion for Embodied Control}
\author{Seonhye Gu\thanks{This work was independently conducted by the author as personal research while affiliated with the AI-Based Surgical Robot Innovation Lab. The affiliation is provided for identification purposes only and does not imply official laboratory output, institutional endorsement, or sponsorship.}\\
\small AI-Based Surgical Robot Innovation Lab}
\date{}
\begin{document}
\maketitle


\section{Protocol Summary}

The confirmatory experiment tests whether a structured target-motion residual that persists after zero-order parameter repair warrants activation of a prepared velocity state. All design choices in this supplement were fixed before the fresh candidate seeds were executed.

Candidate seeds were fixed as 300-339. Static eligibility was evaluated before dynamic treatment, and the first 10 eligible seeds in numeric order were selected: 300, 301, 302, 303, 304, 305, 307, 308, 310, and 311. Thirty-four of 40 candidates passed. Each selected seed was evaluated in a fresh Isaac process for every arm-condition cell, producing a complete 10-seed by 4-arm by 10-condition dynamic grid plus 20 static-retention cells. Dynamic outcomes did not influence seed inclusion.

\section{Exact Model Definitions}

Let \(p_t\) be the observed target command and \(H=10\) the locked prediction horizon. The zero-order model is

\begin{align*}
\hat p_t &= \alpha p_t + (1-\alpha)\hat p_{t-1}, \\
\hat p_{t+H} &= \hat p_t.
\end{align*}

Episode 1 searches \(\alpha \in \{0.25,0.50,0.75,1.00\}\). The selected repair is \(\alpha=1.00\). This model can reduce current-position lag but cannot produce a nonzero open-loop displacement.

The order-one alternative adds

\begin{align*}
\hat v_t &= \beta\left(p_t-p_{t-1}\right)
  +(1-\beta)\hat v_{t-1}, \\
\hat p_{t+H} &= \hat p_t + H\hat v_t.
\end{align*}

after gate activation. The selected coefficients are \(\alpha=\beta=1.00\). Before activation, arm C returns the same position-only forecast as arm B. The oracle uses the true injected condition velocity and is never included in the primary contrast.

\section{Gate Definition And Controls}

The gate uses an eight-step window and requires at least four deltas. It fires if all of the following are true:

\begin{table}[htbp]
\centering
\small
\begin{tabularx}{\textwidth}{@{}>{\raggedright\arraybackslash}X >{\raggedright\arraybackslash}X@{}}
\toprule
Quantity & Locked threshold \\
\midrule
Mean target speed & \(\ge 0.5\ \mathrm{mm/step}\) \\
Active-delta fraction & \(\ge 0.75\) \\
Directional consistency & \(\ge 0.90\) \\
Constant-velocity fit improvement & \(\ge 0.50\) \\
\bottomrule
\end{tabularx}
\end{table}

Directional consistency is the norm of the summed displacement divided by the sum of displacement norms. Constant-velocity fit improvement compares the RMSE of target deltas under zero velocity with the RMSE of delta differences under constant velocity.

Controls use the same origin and number of steps as the persistent-drift trace:

\begin{table}[htbp]
\centering
\small
\begin{tabularx}{\textwidth}{@{}>{\raggedright\arraybackslash}X >{\raggedright\arraybackslash}X >{\raggedright\arraybackslash}X@{}}
\toprule
Control & Construction & Intended rejection \\
\midrule
M0 static & Repeated fixed target & No motion evidence \\
M1 persistent drift & Fixed nonzero delta each step & Positive structural evidence \\
N1 observation noise & IID Gaussian noise, \(\sigma=0.2\ \mathrm{mm}\) & High-frequency residual without direction \\
N2 single impulse & One eight-delta jump followed by rest & Transient event without persistence \\
\bottomrule
\end{tabularx}
\end{table}

\section{Confirmatory Drift Conditions}

\begin{table}[htbp]
\centering
\small
\caption{Fresh Episode 2 condition set. Diagonal components were selected so that the vector norm matches the stated speed up to rounding.}
\label{tab:s1}
\begin{tabularx}{\textwidth}{@{}>{\raggedright\arraybackslash}X >{\raggedright\arraybackslash}X >{\raggedright\arraybackslash}X >{\raggedright\arraybackslash}X >{\raggedright\arraybackslash}X@{}}
\toprule
ID & dx (mm/step) & dy (mm/step) & Delay & Duration \\
\midrule
C01 & +1.700 & 0.000 & 1 & 23 \\
C02 & +1.900 & 0.000 & 4 & 21 \\
C03 & -1.800 & 0.000 & 1 & 23 \\
C04 & -2.000 & 0.000 & 4 & 20 \\
C05 & 0.000 & +1.700 & 2 & 23 \\
C06 & 0.000 & -1.900 & 4 & 21 \\
C07 & +1.202 & +1.202 & 1 & 23 \\
C08 & -1.273 & +1.273 & 4 & 22 \\
C09 & +1.344 & -1.344 & 2 & 21 \\
C10 & -1.414 & -1.414 & 5 & 20 \\
\bottomrule
\end{tabularx}
\end{table}

The set contains both signs of each axis and all four planar diagonal quadrants. It varies speed, delay, and duration and contains no exact condition from the excluded process-isolated pilot.

\section{Execution And Branch Replay}

Each seed-arm-condition cell starts a new Isaac process. The procedure is:

\begin{enumerate}[leftmargin=*]

\item Set the simulator seed and create the environment.

\item Run the static prefix until the target is within 20 mm for five stable steps, or stop at 200 prefix steps.

\item Record the reset-state and branch-state fingerprints.

\item Replay the shared Episode 1 evidence and fit the arm-specific model using only the locked candidates.

\item Restore the matched branch command and run one Episode 2 condition.

\item At every step, provide the model forecast to the shared controller, restore the true command for evaluation, and record target, end-effector, action, model state, gate state, and terminal information.

\item Close the Isaac process before starting the next cell.

\end{enumerate}

Fresh-process execution was introduced after an excluded pilot demonstrated that sequential conditions inside one process could perturb later reset states. The confirmatory process-isolation rule therefore treats each cell as a separate simulator lifecycle rather than relying on nominal reset determinism.

\section{Validity Checklist}

\begin{table}[htbp]
\centering
\small
\caption{Preregistered execution-validity fields and observed values.}
\label{tab:s2}
\begin{tabularx}{\textwidth}{@{}>{\raggedright\arraybackslash}X >{\raggedright\arraybackslash}X >{\raggedright\arraybackslash}X@{}}
\toprule
Field & Requirement & Observed \\
\midrule
Candidate selection & Fixed first-eligible order & Pass \\
Selected seed count & 10 & 10 \\
Dynamic cells & 400 complete & 400/400 \\
Static-retention cells & 20 complete & 20/20 \\
Maximum branch-start EE gap & \(\le 1\ \mathrm{mm}\) & 0 mm \\
Maximum branch-start command gap & \(\le 0.001\ \mathrm{mm}\) & 0 mm \\
Maximum branch-start target distance & \(\le 21\ \mathrm{mm}\) & 15.396 mm \\
Prefix readiness & All selected cells & Pass \\
\(H=10\) windows & Present in every dynamic cell & Pass \\
Drift exposure & Complete in every dynamic cell & Pass \\
Unexpected environment resets & 0 & 0 \\
Forbidden-region violations & 0 & 0 \\
\bottomrule
\end{tabularx}
\end{table}

No record was excluded after selection. The aggregate artifact sets \texttt{\detokenize{valid_confirmatory}} to true.

\section{Full Confirmatory Outcomes}

\begin{table}[htbp]
\centering
\scriptsize
\caption{Arm-level outcomes over the complete dynamic grid.}
\label{tab:s3}
\begin{tabularx}{\textwidth}{@{}>{\raggedright\arraybackslash}X >{\raggedright\arraybackslash}X >{\raggedright\arraybackslash}X >{\raggedright\arraybackslash}X >{\raggedright\arraybackslash}X >{\raggedright\arraybackslash}X >{\raggedright\arraybackslash}X@{}}
\toprule
Arm & n & Success & Prediction error & Final distance & Forbidden & Resets \\
\midrule
A frozen zero order & 100 & 39\% & 20.081 mm & 20.388 mm & 0 & 0 \\
B repaired zero order & 100 & 76\% & 18.401 mm & 18.905 mm & 0 & 0 \\
C gated constant velocity & 100 & 100\% & 7.595 mm & 5.601 mm & 0 & 0 \\
D oracle velocity & 100 & 100\% & 0.000 mm & 4.087 mm & 0 & 0 \\
\bottomrule
\end{tabularx}
\end{table}

\begin{table}[htbp]
\centering
\small
\caption{Preregistered C-minus-B contrasts.}
\label{tab:s4}
\begin{tabularx}{\textwidth}{@{}>{\raggedright\arraybackslash}X >{\raggedright\arraybackslash}X >{\raggedright\arraybackslash}X >{\raggedright\arraybackslash}X >{\raggedright\arraybackslash}X@{}}
\toprule
Endpoint & Contrast & Crossed-bootstrap 95\% CI & Favorable pairs & Decision \\
\midrule
\(H=10\) prediction error & -10.806 mm & [-11.360, -10.331] mm & 100/100 & H1 pass \\
Fixed-horizon final distance & -13.304 mm & [-13.599, -12.982] mm & 100/100 & H2 pass \\
20 mm resolution rate & +0.24 & [+0.08, +0.45] & 24/100 strict gains & Secondary \\
\bottomrule
\end{tabularx}
\end{table}

Success is thresholded at 20 mm. If both B and C succeed, the pair is tied on the binary endpoint even when C has a smaller continuous final distance. This is why 24/100 pairs show a strict binary gain while 100/100 favor C on final distance.

\section{Condition-Level Effects}

\begin{table}[htbp]
\centering
\small
\caption{C-minus-B effects and success counts by condition.}
\label{tab:s5}
\begin{tabularx}{\textwidth}{@{}>{\raggedright\arraybackslash}X >{\raggedright\arraybackslash}X >{\raggedright\arraybackslash}X >{\raggedright\arraybackslash}X >{\raggedright\arraybackslash}X@{}}
\toprule
Condition & Prediction difference & Final-distance difference & B success & C success \\
\midrule
C01 & -11.769 mm & -12.865 mm & 9/10 & 10/10 \\
C02 & -10.364 mm & -13.166 mm & 9/10 & 10/10 \\
C03 & -12.462 mm & -13.807 mm & 9/10 & 10/10 \\
C04 & -10.000 mm & -13.586 mm & 5/10 & 10/10 \\
C05 & -10.461 mm & -12.801 mm & 10/10 & 10/10 \\
C06 & -10.364 mm & -13.593 mm & 5/10 & 10/10 \\
C07 & -11.768 mm & -12.716 mm & 10/10 & 10/10 \\
C08 & -10.502 mm & -13.192 mm & 9/10 & 10/10 \\
C09 & -10.367 mm & -13.427 mm & 5/10 & 10/10 \\
C10 & -9.998 mm & -13.888 mm & 5/10 & 10/10 \\
\bottomrule
\end{tabularx}
\end{table}

Every condition-level mean clears the favorable -5 mm threshold on both continuous endpoints. These rows are descriptive; confirmatory inference uses the crossed seed-condition bootstrap.

\section{Retention And Gate Specificity}

\begin{table}[htbp]
\centering
\scriptsize
\caption{Static retention and gate controls.}
\label{tab:s6}
\begin{tabularx}{\textwidth}{@{}>{\raggedright\arraybackslash}X >{\raggedright\arraybackslash}X >{\raggedright\arraybackslash}X >{\raggedright\arraybackslash}X >{\raggedright\arraybackslash}X >{\raggedright\arraybackslash}X@{}}
\toprule
Analysis & Group & n & Rate & Wilson 95\% CI & Mean final distance \\
\midrule
Static retention & B repaired zero order & 10 & 100\% & - & 1.501 mm \\
Static retention & C gated constant velocity & 10 & 100\% & - & 1.501 mm \\
Gate & M1 persistent drift & 100 & 100\% & [96.30\%, 100\%] & - \\
Gate & M0 static & 100 & 0\% & [0\%, 3.70\%] & - \\
Gate & N1 observation noise & 100 & 0\% & [0\%, 3.70\%] & - \\
Gate & N2 single impulse & 100 & 0\% & [0\%, 3.70\%] & - \\
\bottomrule
\end{tabularx}
\end{table}

The B and C static trajectories are identical over their overlap. The gate does not activate on the static target, so C reduces exactly to repaired zero order.

\section{Statistical Procedure}

For each continuous endpoint, form a 10 by 10 matrix of paired C-minus-B differences indexed by selected seed and fresh condition. For bootstrap replicate b, independently sample 10 seed indices and 10 condition indices with replacement and compute

\begin{align*}
\bar{\Delta}^{*(b)}
  &= \frac{1}{100}\sum_{a=1}^{10}\sum_{c=1}^{10}
     \Delta_{I_a^{(b)},J_c^{(b)}}.
\end{align*}

The reported interval is the 2.5th and 97.5th percentile of 10,000 stored means. The RNG seed is 20260730. This preserves both crossed sampling factors rather than treating 100 cells as independent. The success-rate contrast is resampled identically but remains secondary. Wilson score intervals summarize individual gate rates.

\section{Reproducibility Environment}

\begin{table}[htbp]
\centering
\small
\begin{tabularx}{\textwidth}{@{}>{\raggedright\arraybackslash}X >{\raggedright\arraybackslash}X@{}}
\toprule
Item & Setting \\
\midrule
Simulator & Isaac Sim 4.1 \\
Framework & Isaac Lab 1.0.0 and ORBIT-Surgical \\
Task & \nolinkurl{Isaac-Reach-Dual-STAR-IK-Rel-Play-v0} \\
Fabric & Disabled \\
Preregistration tag & \nolinkurl{paper002-model-order-confirmatory-v1.0} \\
Bootstrap draws / seed & 10,000 / 20260730 \\
\bottomrule
\end{tabularx}
\end{table}

\section{Simulator Warnings And Scope}

The ORBIT-Surgical assets emitted PhysX warnings for rigid bodies with negative mass or placeholder inertia tensors. Isaac Sim substituted approximate inertia. The warning was stable across arms and conditions and did not cause missing cells, unexpected resets, or branch mismatch. It nevertheless limits physical interpretation. The experiment supports a target-model and controller comparison inside the observed simulator stack; it does not validate force, tissue, mass-property, or hardware behavior.

No human participants, animals, patient data, or clinical procedures were involved. The surgical context names the robot benchmark. It is not evidence of clinical safety or efficacy.

\end{document}
